\selectlanguage{english}
\begin{abstract}
Electric power is essential for all areas of production. Therefore, the condition of the power line distribution equipment is essential for energy quality. However, the power grids extend for miles and even stretches of difficult access, compromising the search for damaged equipment. This article presents a solution for detection and classification of objects in power distribution lines using Convolutional Neural Networks. In the first stage of the project, CNN was trained to detect and classify four types of objects that are part of the power grid. YOLOv3 was the artificial neural network selected to detect and classify the objects. To compose the test and training image bank, a drone was used and image collections were made in 10 different locations. After the training, CNN achieved a IoU of 60.38\%.
\end{abstract}
}
\renewcommand{\abstractname}{Resumo}
\begin{abstract}
 A energia elétrica é algo essencial para todas as áreas de produção. Sendo assim, a condição dos equipamentos é essencial para a distribuição de uma energia de qualidade. Contudo, as redes elétricas se estendem por quilômetros e ainda por trechos de difícil acesso, comprometendo encontrar equipamentos danificados. Neste artigo é apresentado uma solução de detecção e classificação de objetos das linhas de distribuição de energia utilizando Redes Neurais Convolucionais. Na primeira etapa do projeto, a CNN foi treinada para detectar e classificar quatro tipos de objetos que fazem parte da rede elétrica. A rede neural artificial selecionada para detectar e classificar os objetos foi a YOLOv3. Para compor o banco de imagens de teste e treinamento, foi utilizado um drone e efetuadas coletas de imagens em 10 locais distintos. Após o treinamento, a CNN alcançou um IoU de 60,38\%.
\end{abstract}



\renewcommand\IEEEkeywordsname{Palavras-chave}

\hyphenation{op-tical net-works semi-conduc-tor}

\begin{document}

\title{Processos Criativos da Obra-jogo "As Lendas de M'Boi": Uma Experiência de Exploração Interativa entre Mitologia e Arte}

%claudio.mauricio@unioeste.br
\author{\IEEEauthorblockN{Franklin L N Fracaro, Fabiana F F Peres,\\Claudio R M Mauricio}
\IEEEauthorblockA{Centro de Engenharias e Ciencias Exatas\\
Univ Estadual do Oeste do Paraná-Unioeste\\
Foz do Iguaçu, Brasil\\
Email: frafracaro@gmail.com, fabiana.peres@unioeste.br,\\ claudio.mauricio@unioeste.br}
\and
\IEEEauthorblockN{Pablo S Villavicencio, \\Ester Marçal Fér}
\IEEEauthorblockA{Instituto Latino-Americano de Arte, Cultura e História\\
Univ Federal da Integração Latino-Americana-Unila\\
Foz do Iguaçu, Brasil\\
Email: pablo.villavicencio@unila.edu.br, \\ester.fer@unila.edu.br}

}

\maketitle

\begin{abstract}
O presente artigo consiste em uma descrição da obra-jogo Lendas de M’Boi e também uma reflexão sobre o seu processo de criação. A obra-jogo foi concebida por meio de uma rede de estudantes coordenados pelos professores do projeto de extensão JogArt - Laboratório de Interatividade em Jogos, Narrativa e Arte, cujo objetivo central é realizar uma aproximação entre o conceito geral de jogo e obras de arte digital. O artigo possui uma fundamentação que reflete sobre o jogo inserido na cultura, jogos digitais, interatividade em obras de arte digitais,  software e  métodos para o processo de criação.  O objetivo central do artigo é apresentar a obra-jogo de forma reflexiva, explicitando os aspectos audiovisuais e interativos que articula. A construção do espaço do jogo em que as espacialidades virtuais são integradas por meio de dispositivos móveis que adicionam camadas ao espaço físico a ser explorado pelo jogador. Portanto, o jogo faz uso do conceito de palimpsesto, em que o jogador ao mesmo tempo que navega e interage com a máquina, se aprofunda nas histórias imaginárias presentes em lugares específicos da região da tríplice fronteira Brasil, Argentina e Paraguai. As lendas, narradas em áudio, são conectadas a esses lugares e disparadas por meio da interatividade com os elementos visuais do jogo. Assim, associou-se no artigo mitologia indígena, reflexões sobre a construção de espaços jogáveis e narrativas que são importantes para a cultura e identidade de Foz do Iguaçu.
\end{abstract}

\begin{IEEEkeywords}
Desenvolvimento de jogos; Lendas indígenas; Processo criativo.
\end{IEEEkeywords}